\section{Implementation and Code Quality}
\subsection{Code Listing}
Link source code: \href{https://github.com}{Click here}
\subsection{Short Inspection Summary}

\begin{table}[H]
\centering
\begin{tabularx}{\linewidth}{|l|X|}
\hline
\textbf{Tiêu chí} & \textbf{Mô tả} \\
\hline
Đặt tên rõ ràng & Các lớp (Event, EventManager), hàm (create\_event, delete\_event) và biến (event\_date, reminder\_minutes) được đặt tên có ý nghĩa, thể hiện rõ mục đích sử dụng. \\
\hline
Thụt lề nhất quán & Tuân theo chuẩn PEP 8: thụt lề 4 khoảng trắng; sử dụng dòng trống để phân tách các khối logic. \\
\hline
Chú thích và docstring & Mỗi lớp và các hàm công khai đều có docstring; comment được dùng để giải thích các logic không rõ ràng. \\
\hline
Nguyên lý trách nhiệm đơn & Lớp Event chỉ chứa dữ liệu; EventManager xử lý toàn bộ logic CRUD và kiểm tra dữ liệu. \\
\hline
Kiểm tra đầu vào & Sử dụng các hàm private riêng (\_validate\_title, \_validate\_datetime, \_validate\_reminder) giúp hàm create\_event rõ ràng hơn. \\
\hline
\end{tabularx}
\caption{Đánh giá chất lượng mã nguồn}
\end{table}

\begin{table}[H]
\centering
\begin{tabularx}{\linewidth}{|l|l|X|X|}
\hline
\textbf{ID} & \textbf{Loại vấn đề} & \textbf{Vấn đề phát hiện} & \textbf{Cách khắc phục đề xuất} \\
\hline
I-01 & Naming & Dict \_events dùng tên quá chung (e dict) & Đổi thành \_event\_store hoặc \_event\_registry để rõ nghĩa hơn \\
\hline
I-02 & Missing check & create\_event không dùng strip() cho description trước khi lưu & Thêm description = description.strip() trước khi tạo Event \\
\hline
I-03 & Readability & to\_dict() không bao gồm trường created\_at, không nhất quán với constructor & Thêm 'created\_at': self.created\_at vào dict trả về \\
\hline
I-04 & Missing check & reminder\_minutes không có giới hạn trên; giá trị rất lớn vẫn được chấp nhận & Thêm kiểm tra giới hạn trên (ví dụ: <= 10080 cho tối đa 1 tuần) \\
\hline
I-05 & Documentation & list\_events() thiếu docstring mô tả kiểu trả về & Thêm annotation kiểu trả về List[Event] và bổ sung docstring \\
\hline
\end{tabularx}
\caption{Kết quả peer-review}
\end{table}

\subsection{Code Smells and Refactoring}
\subsubsection{Code Smell 1 - Long Method (create\_event)}
\noindent\textbf{Vấn đề:} \\
Hàm \texttt{create\_event} gọi ba hàm kiểm tra (validation) riêng biệt. Dù mỗi hàm đều nhỏ, nhưng khi hệ thống mở rộng, nhiều kiểm tra khác (ví dụ: kiểm tra sự kiện trùng lặp, kiểm tra định dạng địa điểm) sẽ tiếp tục được thêm trực tiếp vào \texttt{create\_event}, khiến hàm này trở nên dài và khó quản lý.

\vspace{5pt}

\noindent\textbf{Refactoring:} \\
Tách ra một hàm duy nhất \texttt{\_validate\_all(title, date, time, reminder)} để gọi tất cả các hàm kiểm tra con. Khi đó, \texttt{create\_event} chỉ cần gọi \texttt{\_validate\_all(...)} để giữ độ dài hàm dưới 10 dòng. Cách này tuân theo mẫu refactoring \textit{Extract Method}.

\begin{lstlisting}[style=python]
# Before (smell)
def create_event(self, title, date, time, ...):
    self._validate_title(title)
    self._validate_datetime(date, time)
    self._validate_reminder(reminder_minutes)
    # 10+ more lines 

# After (refactored)
def _validate_all(self, title, date, time, reminder):
    self._validate_title(title)
    self._validate_datetime(date, time)
    self._validate_reminder(reminder)

def create_event(self, title, date, time, ...):
    self._validate_all(title, date, time, reminder_minutes)
    new_event = Event(...)
    self._events[new_event.event_id] = new_event
    return new_event
\end{lstlisting}

\subsubsection{Code Smell 2 - Duplicate Logic (get\_event reuse)}
\noindent\textbf{Vấn đề:} \\
Cả \texttt{delete\_event} và các hàm trong tương lai như \texttt{update\_event} đều cần lấy sự kiện theo ID. Nếu logic tìm kiếm (kiểm tra tồn tại key, ném lỗi \texttt{KeyError}) được viết trực tiếp trong từng hàm, nó sẽ bị lặp lại.

\vspace{6pt}

\noindent\textbf{Refactoring:} \\
Hàm \texttt{get\_event} hiện tại đã gom toàn bộ logic này. Cách cải thiện là đảm bảo mọi hàm cần truy cập sự kiện đều gọi \texttt{get\_event(event\_id)} thay vì truy cập trực tiếp \texttt{self.\_events[event\_id]}. Điều này tuân theo nguyên lý \textit{DRY (Don't Repeat Yourself)}.

\begin{lstlisting}[style=python]
# Before (smell) - direct dict access in delete_event
def delete_event(self, event_id):
    if event_id not in self._events:          # duplicated check
        raise KeyError(...)
    del self._events[event_id]

# After (refactored) - reuse get_event
def delete_event(self, event_id):
    event = self.get_event(event_id)          # centralized check
    del self._events[event_id]
    return True
\end{lstlisting}
